% qualification-checklist.tex
% Document 7 - Center Qualification and Handoff Checklist
%
% INTERNAL USE ONLY. This document is completed by center leadership and
% engineering staff to decide whether to proceed, clarify further, or
% decline a request. It is never distributed to the requester -- keep it
% out of any package sent externally.
%
% Build (from this folder, requires XeLaTeX -- see README.md):
%   xelatex -interaction=nonstopmode qualification-checklist.tex
%   xelatex -interaction=nonstopmode qualification-checklist.tex   (rerun for TOC/page refs)
%
% See docs/Engineering_Project_Intake_Plan.md for the full workflow.

\documentclass[11pt]{article}

\usepackage{../../shared/intake-style}
\usepackage{../../shared/form-components}

\SetDocIdentity{Center Qualification and Handoff Checklist}{ENG-FORM-07}{v0.1-draft}
\SetFieldPrefix{qualchk}

% Replace these placeholders once branding and markings are confirmed
% (see "Remaining Organization-Specific Decisions" in the intake plan).
% The marking text below is intentionally distinct from the requester-
% facing documents to make internal-only handling unmistakable.
\SetCompanyName{[Engineering Center Name]}
\SetMarkingText{INTERNAL USE ONLY -- DO NOT DISTRIBUTE TO REQUESTER}

\SetDocPurpose{%
INTERNAL DOCUMENT -- DO NOT SEND TO THE REQUESTER. This checklist is for
center leadership and engineering staff only. Its purpose is to record an
honest internal assessment of whether the center has the right expertise
and capacity for a request, how confident the team is in its estimate,
what is still missing, and what the funding/charge-code route will be,
before deciding to proceed, clarify further, or decline. Keep this
document, and anything written on it, separate from the requester-facing
document package for this request.%
}

\begin{document}

\begin{Form}

\CoverPage

\DocumentControlBlock

\LinkedIntakeBlock

\IntakeSection{Expertise and Capacity Assessment}
\TextArea{ExpertiseFit}{Does the center have the right disciplines and expertise for this request? Notes.}{\linewidth}{1.5cm}\\[8pt]
\TextArea{CapacityFit}{Current capacity / availability to take on this request}{\linewidth}{1.5cm}

\IntakeSection{Missing Information}
\Guidance{Information still needed before the team can commit to a
confident estimate.}\\[4pt]
\OpenItemsTable{QualOpen}

\IntakeSection{Estimate Confidence}
\DropDown{EstConfidence}{Estimate confidence}{\linewidth}{Low,Medium,High}\\[8pt]
\TextArea{EstNotes}{Notes on estimate confidence (what would increase it)}{\linewidth}{1.5cm}

\IntakeSection{Funding and Charge-Code Route}
\TextArea{FundingRoute}{Expected funding source and charge-code authorization route}{\linewidth}{1.5cm}

\IntakeSection{Risk Assessment}
\RiskRegisterTable{QualRisk}

\IntakeSection{Decision}
\DropDown{Decision}{Decision}{\linewidth}{Proceed,Clarify Further,Decline}\\[8pt]
\TextArea{DecisionRationale}{Rationale for this decision}{\linewidth}{1.5cm}\\[8pt]
\begin{tabularx}{\textwidth}{XX}
  \TextLine{DecisionMaker}{Decision made by}{\linewidth} &
  \TextLine{DecisionDate}{Date}{\linewidth} \\
\end{tabularx}

\end{Form}

\end{document}
